
%(BEGIN_QUESTION)
% Copyright 2015, Tony R. Kuphaldt, released under the Creative Commons Attribution License (v 1.0)
% This means you may do almost anything with this work of mine, so long as you give me proper credit

Les utvalgte deler av instruksjonsmanualen for ``SEL-RS Rotary Switch'' (dokument SEL-RS Rotary Switch Family Instruction Manual, mai 2013), og svar på følgende spørsmål:

\vskip 10pt

En typisk SEL-RS roterende håndbryter består av ulike moduler stablet på en felles roterende aksel. Tabell 1 på side 5 viser fire ulike modultyper. Oppgi hva modulene heter og forklar hva hver av dem gjør:

\begin{itemize}
\item{} 
\item{} 
\item{} 
\item{} 
\end{itemize}

\vskip 10pt

Forklar hva ``A''-terminalene brukes til i denne bryteren (f.eks. A1+, A1$-$, A2+, A2$-$, osv.).

\vskip 10pt

Figur 13 på side 8 viser vanlige koblinger for valgfri {\it tripping module}. I figuren til høyre er et par LED-er koblet sammen med utløserspolen. Angi funksjonen til LED1 og LED3 -- under hvilke betingelser lyser de?

\vskip 10pt

SEL selger tre ulike typer kontaktmoduler for SEL-RS86 roterende bryter (``lockout''-applikasjoner). Oppgi modultypene og forklar hvor hver type kan brukes.

\vskip 10pt

SEL-RS roterende brytere brukes vanligvis i tre hovedapplikasjoner, gjenspeilet i modellnumrene etter ``RS'': SEL-RS86, SEL-RS52 og SEL-RS43. Oppgi applikasjonene og hva modellnumrene refererer til.

\vskip 20pt \vbox{\hrule \hbox{\strut \vrule{} {\bf Forslag til sokratisk diskusjon} \vrule} \hrule}

\begin{itemize}
\item{} Denne typen panelmontert bryter kalles {\it modulær} eller {\it stakkbar}. Forklar hva begrepene betyr, og hvorfor denne konstruksjonen er nyttig i industrielle styreapplikasjoner.
\item{} En funksjon som finnes på enkelte modeller er en {\it target}. Forklar hva ``target'' er og hva den kan brukes til i et elektrisk vernesystem.
\end{itemize}

\underbar{file i00825}
%(END_QUESTION)





%(BEGIN_ANSWER)

\noindent
{\bf Delsvar:}

\vskip 10pt

Modellnumrene refererer til ANSI/LIII-koden for applikasjonen:

\begin{itemize}
\item{} SEL-RS86 = hjelpelåsrelé (auxiliary lockout relay)
\item{} SEL-RS52 = manuell åpne/lukke for effektbryter
\item{} SEL-RS43 = manuell overføring eller velgerfunksjon
\end{itemize}


%(END_ANSWER)





%(BEGIN_NOTES)

Tilgjengelige modultyper for SEL-RS-familien:

\begin{itemize}
\item{} LED-modul (frontpanel-lys)
\item{} Trip-modul (utløsersolenoid, for 86 lockout-applikasjoner)
\item{} Mekanismemodul (fronthåndtak)
\item{} Kontaktmodul (bryterkontakter påvirket av roterakslingen)
\end{itemize}

\vskip 10pt

``A''-terminalene kobles til anode- og katodesiden på hver frontpanel-LED i LED-modulen.

\vskip 10pt

LED1 lyser når NC-kontakten er lukket (bryteren i ``reset''-stilling) og det er kontinuitet gjennom utløserspolen. LED3 lyser når en enhet gir et ``trip''-signal til bryterens utløsersolenoid.

\vskip 10pt

De tre kontaktmodultypene for SEL-RS86 har alle fire kontakter hver, men skiller seg i ``normal'' kontaktstatus. Én type har to NC og to NO kontakter. En annen har fire NO kontakter. Den siste har fire NC kontakter.

\vskip 10pt

Modellnumrene refererer til ANSI/LIII-koden for applikasjonen:

\begin{itemize}
\item{} SEL-RS86 = hjelpelåsrelé (auxiliary lockout relay)
\item{} SEL-RS52 = manuell åpne/lukke for effektbryter
\item{} SEL-RS43 = manuell overføring eller velgerfunksjon
\end{itemize}


%INDEX% Reading assignment: SEL model RS rotary hand switch

%(END_NOTES)
